% =============================================================================
%  Notas de clase — Sesión 11: Funciones Recursivas
%  Programación Avanzada — Universidad Panamericana
%  Compilar: pdflatex sesion11_funciones_recursivas.tex
% =============================================================================
\documentclass[12pt, oneside]{article}
\usepackage{progavanzada}
\usepackage{multicol}

\begin{document}

\portada{Sesión 11}{Funciones Recursivas}{2026}

\tableofcontents
\newpage

% =============================================================================
\section{La recursión antes de las computadoras}
% =============================================================================

Mucho antes de que existiera C++, o cualquier lenguaje de programación,
los matemáticos ya llevaban siglos usando la recursión sin llamarla así.

En 1202, el matemático italiano \textbf{Leonardo de Pisa} —conocido como
Fibonacci— publicó su \textit{Liber Abaci} con el siguiente problema:

\begin{quote}
  \textit{``Un par de conejos produce otro par cada mes.
  Los conejos tardan un mes en madurar y, a partir del segundo mes,
  producen un par nuevo cada mes.
  ¿Cuántos pares habrá después de un año?''}
\end{quote}

La solución es la célebre \textbf{sucesión de Fibonacci}:
$1, 1, 2, 3, 5, 8, 13, 21, 34, 55, \ldots$
donde cada número es la suma de los dos anteriores.

Escrito formalmente:
\[
  F(1) = 1 \qquad F(2) = 1 \qquad F(n) = F(n-1) + F(n-2)
\]

Eso \textit{es} una definición recursiva: un valor definido en términos de
versiones anteriores de sí mismo, con casos base que detienen la cadena.
Los matemáticos usan este estilo desde hace siglos; nosotros simplemente
lo vamos a traducir a código.

\begin{tecnico}[Gödel y las funciones recursivas primitivas]
  En 1931 el matemático \textbf{Kurt Gödel} usó lo que hoy llamamos
  \textit{funciones recursivas primitivas} en su célebre teorema de
  incompletitud para codificar demostraciones matemáticas como números.
  La idea de que una función puede definirse en términos de sí misma con
  un rango más pequeño fue formalizada por Gödel mucho antes de que hubiera
  procesadores capaces de ejecutarla.\\
  \textit{Referencia:} Gödel, K. \textit{Über formal unentscheidbare Sätze
  der Principia Mathematica}. Monatshefte für Mathematik, 1931.
\end{tecnico}

\begin{tecnico}[La primera implementación: LISP (1958)]
  El primer lenguaje de programación que soportó funciones recursivas
  de forma nativa fue \textbf{LISP}, diseñado por \textbf{John McCarthy}
  en el MIT en 1958.  McCarthy quería un lenguaje capaz de procesar listas
  de manera similar a como los matemáticos definen funciones: en términos
  de sí mismas.  Antes de LISP, los lenguajes como FORTRAN no tenían una
  pila de llamadas real y la recursión \textit{simplemente no era posible}.\\
  \textit{Referencia:} McCarthy, J. \textit{Recursive Functions of Symbolic
  Expressions and Their Computation by Machine}. CACM, 1960.
\end{tecnico}

\newpage

% =============================================================================
\section{¿Qué es una función recursiva?}
% =============================================================================

Una \textbf{función recursiva} es una función que, para resolver un problema,
se llama a sí misma con una versión \textit{más pequeña} del mismo problema.

La clave es esa palabra: \textbf{más pequeña}.  No puede llamarse con el mismo
problema (recursión infinita) ni con uno más grande (divergencia).
Cada llamada debe acercar el problema a un punto donde la respuesta sea
inmediata y no requiera más llamadas.

\begin{center}
\begin{tabular}{ll}
  \toprule
  \textbf{Parte} & \textbf{Descripción} \\
  \midrule
  \textbf{Caso base}      & La versión del problema tan pequeña que se resuelve
                            directamente, sin más llamadas. \\[4pt]
  \textbf{Caso recursivo} & Reduce el problema un paso y delega el resto a
                            una llamada con un argumento más pequeño. \\
  \bottomrule
\end{tabular}
\end{center}

\begin{advertencia}
  Si falta el \textbf{caso base}, la función se llama infinitamente y el
  programa termina con un \textbf{stack overflow} (desbordamiento de pila).
  Si el caso recursivo \textbf{no reduce} el problema, pasa lo mismo.
  Ambos errores son los más comunes al aprender recursión.
\end{advertencia}

% =============================================================================
\section{El primer ejemplo: suma acumulada}
% =============================================================================

Queremos calcular $1 + 2 + 3 + \cdots + n$ de forma recursiva.
La observación clave es:
\[
  \text{suma}(n) = n + \underbrace{(1 + 2 + \cdots + (n-1))}_{\text{suma}(n-1)}
\]

Con eso, la definición completa es:

\begin{center}
\begin{tabular}{ll}
  \toprule
  \textbf{Caso} & \textbf{Definición} \\
  \midrule
  $n = 0$ (caso base)      & $\text{suma}(0) = 0$ \\
  $n > 0$ (caso recursivo) & $\text{suma}(n) = n + \text{suma}(n-1)$ \\
  \bottomrule
\end{tabular}
\end{center}

\subsection{Traza de \texttt{suma(4)}}

\begin{ejemplo}[Desenrollar la llamada]
\begin{center}
\begin{tabular}{ll}
  \toprule
  \textbf{Llamada} & \textbf{Expansión} \\
  \midrule
  \texttt{suma(4)} & $= 4 + \texttt{suma(3)}$ \\
  \texttt{suma(3)} & $= 3 + \texttt{suma(2)}$ \\
  \texttt{suma(2)} & $= 2 + \texttt{suma(1)}$ \\
  \texttt{suma(1)} & $= 1 + \texttt{suma(0)}$ \\
  \texttt{suma(0)} & $= 0$ \quad\textit{(caso base)} \\
  \midrule
  \textit{Regresando:} & \\
  \texttt{suma(1)} & $= 1 + 0 = 1$ \\
  \texttt{suma(2)} & $= 2 + 1 = 3$ \\
  \texttt{suma(3)} & $= 3 + 3 = 6$ \\
  \texttt{suma(4)} & $= 4 + 6 = 10$ \\
  \bottomrule
\end{tabular}
\end{center}
\end{ejemplo}

\begin{consejo}
  Para trazar una función recursiva a mano, dibuja primero todas las
  llamadas ``hacia abajo'' hasta el caso base, y luego sustituye los
  valores ``hacia arriba''.  El árbol tiene forma de V: baja hasta
  el caso base y sube resolviendo.
\end{consejo}

% =============================================================================
\section{El call stack: la memoria detrás de la recursión}
% =============================================================================

Cada vez que una función se llama, el sistema reserva un bloque de memoria
llamado \textbf{stack frame} (marco de pila).  Ese bloque guarda:

\begin{itemize}
  \item Los parámetros recibidos (\texttt{x = 4}, \texttt{x = 3}, \ldots)
  \item Las variables locales
  \item La \textbf{dirección de retorno}: a qué línea regresar cuando la
        función termine
\end{itemize}

Los frames se \textbf{apilan} conforme se hacen llamadas y se
\textbf{desapilan} en orden inverso al regresar.
Este mecanismo es el mismo que viste en la sesión 01 con la analogía
del escritorio: la \textit{RAM} acumula los frames como platos en una pila.

\subsection{La pila durante \texttt{suma(3)}}

\begin{verbatim}
  Estado JUSTO ANTES de que suma(0) retorne:

  ┌─────────────────────────┐  ← tope (llamada más reciente)
  │  suma(0)   x = 0        │  retorna 0 de inmediato
  ├─────────────────────────┤
  │  suma(1)   x = 1        │  espera el resultado de suma(0)
  ├─────────────────────────┤
  │  suma(2)   x = 2        │  espera el resultado de suma(1)
  ├─────────────────────────┤
  │  suma(3)   x = 3        │  espera el resultado de suma(2)
  ├─────────────────────────┤
  │  main()                 │  espera el resultado de suma(3)
  └─────────────────────────┘  ← fondo

  Después, los frames se desapilan:
  suma(0)→0   suma(1)→1   suma(2)→3   suma(3)→6   main imprime 6
\end{verbatim}

\begin{recuerda}
  La profundidad de la recursión (número de frames apilados
  al mismo tiempo) está limitada por el tamaño de la pila del proceso,
  que en la mayoría de sistemas es \textbf{1--8 MB}.
  Para \texttt{suma(n)}, la profundidad máxima es $n + 1$ frames.
  Con $n = 100{,}000$ es posible llegar al límite.
\end{recuerda}

% =============================================================================
\section{Sin caso base: stack overflow}
% =============================================================================

Si se omite el caso base, la pila crece sin detenerse:

\begin{verbatim}
int sumaSinBase(int x) {
    return x + sumaSinBase(x - 1);   // ← nunca para
}
\end{verbatim}

\begin{verbatim}
  sumaSinBase(10)
  sumaSinBase(9)
  sumaSinBase(8)
  ...
  sumaSinBase(-50000)
  ...
  ¡CRASH!  Segmentation fault / Stack overflow
\end{verbatim}

\begin{advertencia}
  El nombre del sitio de preguntas técnicas más famoso del mundo,
  \textbf{Stack Overflow} (\url{https://stackoverflow.com}), es
  exactamente esta situación: la broma interna del fundador
  Jeff Atwood sobre el error más icónico al aprender a programar.
\end{advertencia}

% =============================================================================
\section{Ejemplos adicionales}
% =============================================================================

\subsection{Factorial}

\[
  0! = 1 \qquad n! = n \times (n-1)!
\]

\begin{ejemplo}[Traza de \texttt{factorial(5)}]
\[
  5! = 5 \times 4! = 5 \times 4 \times 3! = 5 \times 4 \times 3 \times 2!
     = 5 \times 4 \times 3 \times 2 \times 1! = 5 \times 4 \times 3 \times 2 \times 1 \times 0!
     = 120
\]
\end{ejemplo}

\begin{consejo}
  El factorial crece extremadamente rápido: $13! > 6 \times 10^9$, que ya
  supera el rango de un \texttt{int} de 32 bits.  Usa \texttt{long long}
  para factoriales mayores que 12.
\end{consejo}

\subsection{Potencias}

\[
  b^0 = 1 \qquad b^n = b \times b^{n-1}
\]

\begin{ejemplo}[Traza de \texttt{pot(2, 5)}]
\[
  2^5 = 2 \times 2^4 = 2 \times 2 \times 2^3 = \cdots = 2^5 = 32
\]
Cada llamada reduce el exponente en 1 hasta llegar a $b^0 = 1$.
\end{ejemplo}

\subsection{Sucesión de Fibonacci}

\[
  F(1) = 1, \quad F(2) = 1, \quad F(n) = F(n-1) + F(n-2)
\]

Fibonacci es uno de los ejemplos más naturales de recursión \textbf{doble}:
el caso recursivo hace \textit{dos} llamadas en lugar de una.

\begin{advertencia}
  La implementación recursiva directa de Fibonacci es \textbf{ineficiente}:
  $F(40)$ requiere más de $10^8$ llamadas porque calcula los mismos valores
  una y otra vez.  Este problema se resuelve con \textit{memoización} o
  programación dinámica — temas de cursos posteriores.
\end{advertencia}

\begin{tecnico}[Fibonacci y la naturaleza]
  La sucesión de Fibonacci aparece en la disposición de las semillas de
  un girasol, las espirales de las piñas de pino y el número de pétalos de
  muchas flores.  La razón $F(n+1)/F(n)$ converge a la \textbf{proporción
  áurea} $\varphi \approx 1.618\ldots$, que también aparece en la
  arquitectura griega y en el arte del Renacimiento.  Lo que comenzó como
  un problema de conejos en 1202 terminó siendo una de las constantes
  matemáticas más observadas en la naturaleza.\\
  \textit{Referencia:} Livio, M. \textit{The Golden Ratio}.
  Broadway Books, 2002.
\end{tecnico}

% =============================================================================
\section{Recursivo vs.\ iterativo}
% =============================================================================

Cualquier función recursiva puede reescribirse con un ciclo, y viceversa.
La elección depende de la claridad y el contexto.

\begin{center}
\begin{tabular}{lll}
  \toprule
  \textbf{Criterio} & \textbf{Recursivo} & \textbf{Iterativo} \\
  \midrule
  Legibilidad      & Alta cuando el problema es naturalmente recursivo
                   & Alta cuando hay un ciclo obvio \\[4pt]
  Memoria          & Un frame por llamada activa (pila)
                   & Solo las variables del ciclo \\[4pt]
  Riesgo de fallo  & Stack overflow si la profundidad es grande
                   & No tiene este riesgo \\[4pt]
  Casos ideales    & Árboles, backtracking, divide y vencerás
                   & Sumas, búsquedas simples, arreglos \\
  \bottomrule
\end{tabular}
\end{center}

\begin{recuerda}
  En este curso veremos que los algoritmos sobre \textbf{árboles} y
  \textbf{grafos} son casi siempre más elegantes de forma recursiva.
  Las \textbf{Torres de Hanoi} de la sesión 12 son el ejemplo clásico donde
  la solución iterativa equivalente es considerablemente más difícil de
  entender.
\end{recuerda}

% =============================================================================
\section{Ejercicios}
% =============================================================================

Resuelve los siguientes ejercicios \textbf{a mano, en papel}.
Para cada función recursiva, muestra la traza completa de llamadas
(caso base y regreso de valores).

\begin{enumerate}

  \item Calcula \texttt{suma(6)} trazando todas las llamadas.
        ¿Cuántos frames se apilan simultáneamente como máximo?

  \item Calcula \texttt{factorial(6)} trazando la expansión y la
        resolución inversa.  ¿Cuántas multiplicaciones se realizan?

  \item Calcula \texttt{pot(3, 4)} paso a paso.
        Verifica que el resultado sea $81$.

  \item La siguiente función tiene un error.  Identifica el problema,
        explica qué pasaría en ejecución y escribe la versión corregida.
        \[
          \texttt{int malo(int n) \{ return n + malo(n + 1); \}}
        \]

  \item Escribe la \textbf{definición recursiva} (en notación matemática,
        no en código) de las siguientes funciones:
        \begin{enumerate}
          \item La suma de los primeros $n$ números impares:
                $1 + 3 + 5 + \cdots + (2n-1)$.
          \item El producto $n \times (n+1) \times (n+2) \times \cdots \times m$
                (con $n \leq m$).
        \end{enumerate}

  \item Fibonacci.  Calcula $F(7)$ trazando el árbol de llamadas.
        Cuenta cuántas veces se calcula $F(3)$ en el proceso.
        ¿Ves el problema de eficiencia?

  \item \textbf{Problema de reflexión.}
        La función \texttt{suma(n)} tiene profundidad de recursión $n+1$.
        Si el sistema admite una pila de 8 MB y cada frame ocupa 32 bytes,
        ¿cuál es la profundidad máxima teórica?
        ¿Qué tan grande puede ser $n$ antes de provocar stack overflow?

\end{enumerate}

% =============================================================================
\section{Resumen}
% =============================================================================

\begin{recuerda}
\begin{itemize}
  \item Una función recursiva se llama a sí misma con un argumento
        \textbf{más pequeño}; cada llamada se acerca al caso base.
  \item Toda función recursiva correcta tiene:
        (1) un \textbf{caso base} que retorna sin llamar, y
        (2) un \textbf{caso recursivo} que reduce el problema.
  \item Cada llamada crea un \textbf{stack frame} en la pila;
        los frames se desapilan al retornar, en orden inverso.
  \item Sin caso base → \textbf{stack overflow}.
        Si el argumento no decrece → también.
  \item La recursión es más elegante que los ciclos en problemas con
        estructura naturalmente recursiva (árboles, particiones, etc.).
  \item Fibonacci recursivo directo es \textbf{exponencialmente ineficiente};
        hay soluciones mejores que veremos más adelante.
\end{itemize}
\end{recuerda}

\end{document}
